% !TEX root = ../main.tex
% Triage pipeline as implemented in apps/server/src/modules/triage/triage.engine.ts
\begin{figure}[t]
\centering
\begin{tikzpicture}[
  node distance=4mm and 3mm,
  box/.style={draw, rounded corners=2pt, align=center, font=\footnotesize, minimum height=7mm, inner sep=3pt},
  rule/.style={box, fill=red!8},
  llm/.style={box, fill=blue!8},
  human/.style={box, fill=green!10},
  arr/.style={-{Stealth[length=2mm]}, thick}
]
\node[box] (in) {Vitals +\\symptoms};
\node[rule, right=of in] (safe) {Deterministic\\safety engine};
\node[llm, right=of safe] (rag) {Guideline\\retrieval};
\node[llm, below=of rag] (llm) {LLM triage\\(structured)};
\node[rule, left=of llm] (rec) {Reconcile:\\never below\\safety verdict};
\node[human, left=of rec] (doc) {Doctor\\review};

\draw[arr] (in) -- (safe);
\draw[arr] (safe) -- (rag);
\draw[arr] (rag) -- (llm);
\draw[arr] (llm) -- (rec);
\draw[arr] (rec) -- (doc);
\draw[arr, dashed] (safe) -- node[right, font=\scriptsize] {verdict} (rec);
\end{tikzpicture}
\caption{JeevanSetu triage pipeline. Red: deterministic components; blue: LLM components; green: human-in-the-loop.}
\label{fig:pipeline}
\end{figure}
